\section{The exact interfaces being combined}\label{sec:interfaces}
The new schedules change how work is allocated across recursive levels.
Their correctness uses the same deformation, recovery, and child-factorization
identities as the supplied preprint and refinement packets. We spell out the
interfaces to make the combination checkable. We do not substitute the
weaker formalized endpoint for the stronger manuscript bound, nor infer a
formal proof from finite tests.

\subsection{Asymmetric packing and the fresh side}
At an active call with $M$ wires, choose a packet length $L=2^a$, an alphabet
$r\ge3$, and $Q=r^a$ whole-vertex labels on each side. Let $R,C$ bound the
row and column degrees. The asymmetric packing theorem gives labels with
maximum ordered-pair load at most
\begin{equation}\label{eq:load}
 \left(1+\frac{2}{h_0-1}\right)
 \left(\frac{M}{Q^2}+\frac{R}{Q}\right)
 +2(h_0C+1)\ln Q,
 \qquad h_0\ge2.
\end{equation}
This is an imported theorem of the asymmetric packet~\cite{asymmetric}.
Its construction is useful here: greedily balance row stub counts, then
assign column labels by exact conditional expectations of
$\sum_{x,y}z^{X_{xy}}$, with $z=1+1/(h_0C)$. The average of the candidate
expectations equals the previous expectation, so a minimizing label never
increases it. Convexity of $z^w$ on $0\le w\le C$ proves
\eqref{eq:load}. No real logarithm is evaluated to make an algorithmic choice.

Write $\lambda=\ceil{\log_2s}$ and take $h_0=\max\{2,\lambda\}$.
Rounding the maximum load upward to a multiple of $L$ yields a common padded
mass $P=(1+o(1))M$ whenever, uniformly over active calls,
\begin{equation}\label{eq:packing-conditions}
 \frac{QR}{M}=o(1),\qquad
 \frac{Q^2C\lambda^2}{M}=o(1),\qquad
 \frac{LQ^2}{M}=o(1).
\end{equation}
The last condition pays for rounding. The orientation construction then
retains exactly $P/r$ padded ports per symbol on each old side. Each child
has at most
\begin{equation}\label{eq:child}
 (2/r+o(1))M
\end{equation}
wires. At a coordinate, exactly a fraction $1-1/r$ of ordered label pairs
have unequal symbols and create crossing modes.

The crucial degree invariant is asymmetric. Each child places all retained
parent tensors on one side and newly formed packet tensors on the other.
The retained tensors acquire no ports, and each new packet tensor has degree
at most $2L$. Thus at depth $d\ge1$,
\begin{equation}\label{eq:degrees}
 R_d\le\ceil{s^{9/10}},\qquad C_d\le2^{a_{d-1}+1},
\end{equation}
provided every preceding packet degree was below $s^{9/10}$. Our proofs check
that premise explicitly. At the root both sides can have degree
$\ceil{s^{9/10}}$.

\subsection{Full crossing-block modes preserve the interface}
For a retained interpolation node of valuation $i<u\le a$, put
$D_i=2^{a-i}$. The joint-Cauchy packet constructs integral, group-local
separated modes for the full crossing orientation~\cite{joint}. Their
certified valuations are
\begin{equation}\label{eq:weights}
 w_q=a+1-i+4q-2\pop(q),\qquad 0\le q<D_i/2.
\end{equation}
Each rank index has $2^i$ copies per orientation and there are two
orientations per unequal-symbol packet. Consequently the number of copies
per rank index is bounded by
\begin{equation}\label{eq:families}
 c_i\le\bigl(2(1-1/r)+o(1)\bigr)M/D_i.
\end{equation}
The all-level decomposition is essential: the older annulus construction's
slopes tend to three. A local slope-four calculation at $i=0$ alone would
not justify replacing three by four uniformly.

For clarity, the full-block construction uses slots
$\sigma(j)=(j\bmod2^i)+2^{i+1}\floor{j/2^i}$ and a finite telescoping
Cauchy identity. At pivot $b=q2^i+t$, both separated factors have divisibility
$2^{E_b}$, where $E_b=q+\nu_2(q!)=2q-\pop(q)$. The crossing scalar contributes
$a+1-i$ more powers of two, giving~\eqref{eq:weights}. A mode reads columns
in one group and supplies coefficients only to the other group. It changes
neither child wires nor the degree invariant~\eqref{eq:degrees}. This is why
it composes with asymmetric packing without a hidden multiplicity penalty.

The retained tuple constraint is $\sum_qw_q\beta_q<p_i$. The exact source
identity expresses the deformed contraction as a sum over such tuples of a
product of $r$ child contractions. Features and coefficients can be computed
in polynomial bit time per elementary operation in the same coefficient
rings as the source. Known factors of two are cancelled before unit inversion.

\subsection{Sharp Hermite precision with a thin gap}
The exact linear-taper count and recovery losses are prior
results~\cite{contraction,precision}. Write
\begin{equation}\label{eq:hermite}
 h=\ceil{\frac{M(2L-1)+1}{L(2u-2+2^{1-u})}},\qquad
 q_i=h(u-i),\qquad
 p_i=p+6hD_i+hu(a+1).
\end{equation}
The total number of Hermite conditions is
$hL(2u-2+2^{1-u})$, greater than the degree $M(2L-1)$ of the deformation.
The displayed $p_i$ is a sufficient working precision; exact smaller local
losses could also be used.

Both new schedules set
\begin{equation}\label{eq:gap}
 g=\ceil{\log_2\lambda}^{2},\qquad
 a_d=\sum_{j=0}^du_j+(d+1)g,
\end{equation}
and stop at $H=\ceil{s/\lambda}$. At $O(\log\lambda)$ active depths,
$a_d,u_d=O(\lambda)$ and the accumulated previous precision is bounded by
$O(s\lambda^2 2^{a_{d-1}})$. Since $i<u_d$,
$D_i\ge2^{a_{d-1}+g+1}$. It follows that
\[
 \frac{p}{MD_i/u_d}=O(\lambda^4 2^{-g})=o(1),\qquad
 \frac{hu_d(a_d+1)}{MD_i/u_d}=o(1).
\]
When $u_d\to\infty$ uniformly, $h=(1+o(1))M/u_d$, and therefore
\begin{equation}\label{eq:precision}
 p_i\le(6+o(1))MD_i/u_d.
\end{equation}
We prove the uniform growth and depth bound for each schedule below.

\subsection{The resulting local cost}
Lowering~\eqref{eq:weights} to $\widehat w_q=1+4q-2\pop(q)$ and extending
the finite rank interval only enlarge the tuple set. The joint packet's
weighted-partition estimate is
\[
 \sum_{q\ge0}-\ln(1-e^{-\theta\widehat w_q})
 =\frac{\pi^2}{24\theta}+O\bigl(\log^2(2/\theta)\bigr).
\]
Hence $N_i$, the number of retained tuples, satisfies
$\ln N_i\le\theta p_i+c_i\pi^2/(24\theta)+
O(c_i\log^2(2/\theta))$. Optimize using the prescribed upper envelopes
in~\eqref{eq:families} and~\eqref{eq:precision}. The optimizing
$\theta\asymp\sqrt{u_d}/D_i$ tends uniformly to zero, and its error is
lower order because $D_i\ge2^g$. This yields
\begin{equation}\label{eq:local-entropy}
 \log_2 N_i\le\bigl(k(r)+o(1)\bigr)\frac{M}{\sqrt{u_d}},
 \qquad k(r)=\frac{\pi\sqrt{2(1-1/r)}}{\ln2}.
\end{equation}
The factor four in the rank weights is already proved by prior work; the
new contribution is the feasible allocation of the $u_d$ and alphabets
under the changing degree capacities.

\begin{remark}[Creation-time representatives are retained]
An existing tensor is evaluated with the precision at which its description
was created. Later calls reuse its canonical integer representatives.
Canonical lifting is not a ring homomorphism. The new schedules and kernel
implementations do not replace that convention or recursively compose entry
evaluators. Shared immutable descriptions preserve the source's oracle and
space accounting.
\end{remark}
